\section{Capítulo 5: Desarrollo}

En este capítulo se presenta la fase de desarrollo del sistema, en la cual se lleva a cabo la implementación de la arquitectura y el diseño previamente definidos en la etapa de análisis. A partir de los requerimientos funcionales, no funcionales y las reglas de negocio establecidas en el capítulo anterior, se detalla la construcción de los distintos módulos que conforman la plataforma.

\subsection{Desarrollo de la base de datos}

El desarrollo de la base de datos presentado en esta sección se basa en el diagrama relacional y decisiones de diseño de la base de datos. Se llevó a cabo la implementación en el sistema gestor de bases de datos SQL Server, traduciendo cada entidad, atributo y relación en estructuras físicas mediante sentencias T-SQL. 

\subsubsection{Creación del esquema}

Durante esta etapa, se implementaron restricciones de integridad y definición de claves primarias y foráneas, con el objetivo de garantizar la consistencia, eficiencia y correcto funcionamiento de la aplicación web.

\textbf{A. Gestión de usuarios}

Este módulo está compuesto por las tablas \texttt{Usuarios}, \texttt{Tutor} y \texttt{Alumno}, las cuales representan la estructura principal de acceso y perfiles del sistema. La tabla \texttt{Estatus} actúa como catálogo compartido y debe crearse previamente, ya que es referenciada por \texttt{Usuarios} mediante una clave foránea.

\begin{lstlisting}[language=SQL, caption={Implementación simplificada de entidades de usuario}]
    CREATE TABLE Usuarios (
        id_usuario INT IDENTITY(1,1) PRIMARY KEY,
        usuario VARCHAR(50) NOT NULL UNIQUE,
        contrasena_cifrada VARCHAR(255) NOT NULL,
        tipo_usuario VARCHAR(30) NOT NULL,
        id_estatus INT NOT NULL,
        FOREIGN KEY (id_estatus)
            REFERENCES Estatus(id_estatus)
    );
    
    CREATE TABLE Alumno (
        id_alumno INT IDENTITY(1,1) PRIMARY KEY,
        grado VARCHAR(15) NOT NULL,
        grupo VARCHAR(15) NOT NULL,
        id_tutor INT NOT NULL,
        FOREIGN KEY (id_tutor)
            REFERENCES Tutor(id_tutor)
    );
    \end{lstlisting}

\textbf{B. Proceso de evaluación}

Este módulo agrupa las tablas centrales del flujo de evaluación: \texttt{Ejercicio\_Tutor}, \texttt{Intento} y \texttt{Analisis\_ortografico}. La tabla \texttt{Ejercicio\_Tutor} materializa la relación M:N entre tutores y ejercicios del catálogo, incorporando atributos propios de la asignación. La tabla \texttt{Intento} registra cada ejecución de un ejercicio por parte del alumno, almacenando la imagen capturada y el texto reconocido por el módulo OCR. La tabla \texttt{Analisis\_ortografico} almacena los resultados del procesamiento NLP, utilizando un campo JSON para representar las sugerencias de corrección de forma flexible.

\begin{lstlisting}[language=SQL, caption={Estructura simplificada del proceso de evaluación}]
    CREATE TABLE Ejercicio_Tutor (
        id_ejercicio_tutor INT IDENTITY(1,1) PRIMARY KEY,
        id_tutor INT NOT NULL,
        id_ejercicio INT NOT NULL,
        fecha_limite DATETIME NULL,
        id_estatus INT NOT NULL
    );
    
    CREATE TABLE Intentos (
        id_intento INT IDENTITY(1,1) PRIMARY KEY,
        id_alumno INT NOT NULL,
        id_ejercicio_tutor INT NOT NULL,
        texto_detectado_ocr VARCHAR(MAX),
        fecha_envio DATETIME2 DEFAULT GETDATE()
    );
\end{lstlisting}

\subsubsection{Implementación de restricciones e integridad referencial}

Con el objetivo de garantizar la consistencia y calidad de los datos almacenados en la aplicación web, se definieron restricciones de integridad a nivel de base de datos. Estas restricciones operan de forma independiente a la lógica del backend, asegurando que ninguna operación pueda comprometer la coherencia del modelo de datos.

A continuación se describen las restricciones implementadas agrupadas por categoría.

\paragraph{Unicidad de campos críticos}

Se establecieron restricciones \texttt{UNIQUE} en los campos que deben identificar de forma exclusiva a cada registro dentro del sistema, evitando duplicados que comprometan la integridad de la información.

\begin{lstlisting}[
    language=SQL,
    caption={Restricciones de unicidad en campos críticos},
    label={lst:restricciones_unicidad}
]
CONSTRAINT UQ_Usuarios_usuario 
    UNIQUE (usuario),

CONSTRAINT UQ_Tutor_correo 
    UNIQUE (correo),

CREATE UNIQUE INDEX UQ_Intentos_original
ON Intentos (id_alumno, id_ejercicio_tutor)
WHERE id_recomendacion IS NULL;

CREATE UNIQUE INDEX UQ_ET_activo
    ON Ejercicio_Tutor (id_tutor, id_ejercicio)
    WHERE id_estatus = 1;

CONSTRAINT UQ_Analisis_Ort_intento 
    UNIQUE (id_intento),

CONSTRAINT UQ_Analisis_Cal_intento 
    UNIQUE (id_intento)
\end{lstlisting}

\paragraph{Validación de rangos en scores}

Los atributos de puntuación en las tablas \texttt{Analisis\_Ortografico} y \texttt{Analisis\_Caligrafico} deben estar acotados en un rango de 0 a 10, representando la escala de evaluación del sistema. Se implementaron restricciones \texttt{CHECK} para garantizar que ningún valor fuera de este rango pueda ser insertado o actualizado.

\begin{lstlisting}[language=SQL, caption={Restricciones de rango en atributos de puntuación}]
CONSTRAINT CK_ortografia_score
    CHECK (ortografia_score IS NULL 
        OR (ortografia_score >= 0 AND ortografia_score <= 10)),

CONSTRAINT CK_alineacion_score
    CHECK (alineacion_score IS NULL 
        OR (alineacion_score >= 0 AND alineacion_score <= 10)),

CONSTRAINT CK_tamano_letra_score
    CHECK (tamano_letra_score IS NULL 
        OR (tamano_letra_score >= 0 AND tamano_letra_score <= 10)),

CONSTRAINT CK_espaciado_score
    CHECK (espaciado_score IS NULL 
        OR (espaciado_score >= 0 AND espaciado_score <= 10)),

CONSTRAINT CK_inclinacion_score
    CHECK (inclinacion_score IS NULL 
        OR (inclinacion_score >= 0 AND inclinacion_score <= 10))
\end{lstlisting}

\paragraph{Integridad referencial}

Se definieron claves foráneas en todas las entidades que mantienen relaciones entre sí, garantizando que no puedan existir registros huérfanos en el sistema. La siguiente tabla resume las relaciones de integridad referencial implementadas.

\begin{longtable}{|p{4.5cm}|p{4cm}|p{4cm}|}
\caption{Relaciones de integridad referencial del sistema}
\label{tab:integridad_referencial} \\

\hline
\textbf{Tabla origen} & \textbf{Tabla destino} & \textbf{Cardinalidad} \\
\hline
\endfirsthead

\hline
\textbf{Tabla origen} & \textbf{Tabla destino} & \textbf{Cardinalidad} \\
\hline
\endhead

Usuarios & Estatus & N:1 \\
\hline
Tutor & Usuarios & 1:1 \\
\hline
Alumno & Usuarios & 1:1 \\
\hline
Alumno & Tutor & N:1 \\
\hline
Ejercicio\_Tutor & Tutor & N:1 \\
\hline
Ejercicio\_Tutor & Ejercicio & N:1 \\
\hline
Ejercicio\_Tutor & Estatus & N:1 \\
\hline
Intentos & Alumno & N:1 \\
\hline
Intentos & Ejercicio\_Tutor & N:1 \\
\hline
Analisis\_Ortografico & Intentos & 1:1 \\
\hline
Sugerencia\_ortografica & Analisis\_Ortografico & N:1 \\
\hline
Analisis\_Caligrafico & Intentos & 1:1 \\
\hline

\end{longtable}

\paragraph{Obligatoriedad de campos}

Se establecieron restricciones \texttt{NOT NULL} en todos los atributos que son indispensables para el funcionamiento del sistema, garantizando que ningún registro pueda ser insertado sin la información mínima requerida. Los campos marcados como \texttt{NULL} corresponden a información opcional o generada de forma asíncrona, como el texto detectado por OCR o los scores de análisis, los cuales se completan una vez procesado el intento.

\subsubsection{Implementación de la vista vw\_historial\_alumno}

La vista \texttt{vw\_historial\_alumno} se creó con el objetivo de consolidar en una única consulta predefinida la información distribuida entre las tablas \texttt{Intentos}, \texttt{Analisis\_Ortografico}, \texttt{Analisis\_Caligrafico}, \texttt{Alumno} y \texttt{Ejercicio}. Su implementación responde a tres necesidades concretas del sistema:

\begin{itemize}
    \item \textbf{Rendimiento:} evita la ejecución repetitiva de consultas con múltiples \texttt{JOIN} cada vez que el tutor consulta el historial de un alumno (RF-09) o el alumno revisa sus ejercicios completados (RF-16).
    \item \textbf{Separación de responsabilidades:} el controlador accede al historial mediante una consulta simple sobre la vista, sin necesidad de conocer la estructura interna de las tablas involucradas, lo cual es coherente con el patrón MVC adoptado.
    \item \textbf{Normalización:} permite mantener el modelo de datos normalizado sin sacrificar la facilidad de consulta, ya que la vista actúa como capa de abstracción sobre las relaciones existentes sin almacenar datos redundantes.
\end{itemize}

Se utilizaron \texttt{LEFT JOIN} en las tablas de análisis para contemplar el caso en que un intento recién enviado aún no haya sido procesado por los módulos OCR y NLP, mostrando \texttt{NULL} en los scores mientras el análisis está pendiente.

\begin{lstlisting}[language=SQL, caption={Implementación de la vista vw\_historial\_alumno}]
CREATE VIEW vw_historial_alumno AS
SELECT
    a.id_alumno,
    u.nombre                        AS nombre_alumno,
    e.id_ejercicio,
    e.titulo                        AS titulo_ejercicio,
    e.tipo                          AS tipo_ejercicio,
    i.id_intento,
    i.fecha_envio,
    i.retroalimentacion,
    CASE
        WHEN i.id_recomendacion IS NULL
        THEN 'Original'
        ELSE 'Reintento'
    END                             AS tipo_intento,
    ao.ortografia_score,
    ao.cantidad_errores,
    ac.alineacion_score,
    ac.tamano_letra_score,
    ac.espaciado_score,
    ac.inclinacion_score
FROM Intentos i
INNER JOIN Alumno a
    ON a.id_alumno = i.id_alumno
INNER JOIN Usuarios u
    ON u.id_usuario = a.id_usuario
INNER JOIN Ejercicios_Tutor et
    ON et.id_ejercicio_tutor = i.id_ejercicio_tutor
INNER JOIN Ejercicio e
    ON e.id_ejercicio = et.id_ejercicio
LEFT JOIN Analisis_Ortografico ao
    ON ao.id_intento = i.id_intento
LEFT JOIN Analisis_Caligrafico ac
    ON ac.id_intento = i.id_intento;
GO
\end{lstlisting}

Una vez creada la vista, el historial completo de un alumno se obtiene desde el controlador mediante una consulta directa, filtrando por \texttt{id\_alumno} y ordenando por fecha de envío descendente, sin necesidad de replicar la lógica de JOIN en cada petición.


\begin{lstlisting}[language=SQL, caption={Consulta del historial de un alumno mediante la vista}]
SELECT *
FROM vw_historial_alumno
WHERE id_alumno = @id_alumno
ORDER BY fecha_envio DESC;
\end{lstlisting}

\subsubsection{Implementación de triggers}

Los triggers implementados tienen como objetivo mantener actualizada la tabla \texttt{Progreso\_Alumno} de forma automática cada vez que se registra un nuevo análisis, sin requerir intervención del controlador. Se definieron dos triggers independientes: uno sobre \texttt{Analisis\_Ortografico} y otro sobre \texttt{Analisis\_Caligrafico}, cada uno responsable de actualizar los promedios correspondientes a su dimensión de evaluación.

Ambos triggers manejan dos escenarios: la creación del primer registro de progreso del alumno y la actualización de un registro existente, garantizando la consistencia de los datos en cualquier situación.

\textbf{Trigger sobre Analisis\_Ortografico}
\vspace{0.5cm}
\begin{lstlisting}[language=SQL, caption={Trigger de actualización del promedio ortográfico}]
CREATE TRIGGER trg_actualizar_progreso_ortografia
ON Analisis_Ortografico
AFTER INSERT
AS
BEGIN
    SET NOCOUNT ON;

    DECLARE @id_alumno      INT;
    DECLARE @nuevo_promedio DECIMAL(5,2);

    SELECT @id_alumno = i.id_alumno
    FROM inserted ins
    INNER JOIN Intentos i ON i.id_intento = ins.id_intento;

    SELECT @nuevo_promedio = AVG(ao.ortografia_score)
    FROM Analisis_Ortografico ao
    INNER JOIN Intentos i ON i.id_intento = ao.id_intento
    WHERE i.id_alumno = @id_alumno
    AND ao.ortografia_score IS NOT NULL;

    IF EXISTS (
        SELECT 1 FROM Progreso_Alumno
        WHERE id_alumno = @id_alumno
    )
    BEGIN
        UPDATE Progreso_Alumno
        SET promedio_ortografia = @nuevo_promedio,
            fecha_modificacion  = GETDATE()
        WHERE id_alumno = @id_alumno;
    END
    ELSE
    BEGIN
        INSERT INTO Progreso_Alumno (
            id_alumno,
            promedio_ortografia,
            promedio_aliniacion,
            promedio_tamano_letra,
            promedio_espaciado,
            prommedio_inclinacion,
            fecha_modificacion
        )
        VALUES (
            @id_alumno,
            @nuevo_promedio,
            NULL, NULL, NULL, NULL,
            GETDATE()
        );
    END
END;
GO
\end{lstlisting}
\vspace{0.5cm}
\textbf{Trigger sobre Analisis\_Caligrafico}
\vspace{0.5cm}

\begin{lstlisting}[language=SQL, caption={Trigger de actualización de promedios caligráficos}]
CREATE TRIGGER trg_actualizar_progreso_caligrafico
ON Analisis_Caligrafico
AFTER INSERT
AS
BEGIN
    SET NOCOUNT ON;

    DECLARE @id_alumno              INT;
    DECLARE @prom_alineacion        DECIMAL(5,2);
    DECLARE @prom_tamano_letra      DECIMAL(5,2);
    DECLARE @prom_espaciado         DECIMAL(5,2);
    DECLARE @prom_inclinacion       DECIMAL(5,2);

    SELECT @id_alumno = i.id_alumno
    FROM inserted ins
    INNER JOIN Intentos i ON i.id_intento = ins.id_intento;

    SELECT
        @prom_alineacion    = AVG(ac.alineacion_score),
        @prom_tamano_letra  = AVG(ac.tamano_letra_score),
        @prom_espaciado     = AVG(ac.espaciado_score),
        @prom_inclinacion   = AVG(ac.inclinacion_score)
    FROM Analisis_Caligrafico ac
    INNER JOIN Intentos i ON i.id_intento = ac.id_intento
    WHERE i.id_alumno = @id_alumno
    AND ac.alineacion_score   IS NOT NULL
    AND ac.tamano_letra_score IS NOT NULL
    AND ac.espaciado_score    IS NOT NULL
    AND ac.inclinacion_score  IS NOT NULL;

    IF EXISTS (
        SELECT 1 FROM Progreso_Alumno
        WHERE id_alumno = @id_alumno
    )
    BEGIN
        UPDATE Progreso_Alumno
        SET promedio_aliniacion     = @prom_alineacion,
            promedio_tamano_letra   = @prom_tamano_letra,
            promedio_espaciado      = @prom_espaciado,
            prommedio_inclinacion   = @prom_inclinacion,
            fecha_modificacion      = GETDATE()
        WHERE id_alumno = @id_alumno;
    END
    ELSE
    BEGIN
        INSERT INTO Progreso_Alumno (
            id_alumno,
            promedio_ortografia,
            promedio_aliniacion,
            promedio_tamano_letra,
            promedio_espaciado,
            prommedio_inclinacion,
            fecha_modificacion
        )
        VALUES (
            @id_alumno,
            NULL,
            @prom_alineacion,
            @prom_tamano_letra,
            @prom_espaciado,
            @prom_inclinacion,
            GETDATE()
        );
    END
END;
GO
\end{lstlisting}

\subsection{Desarrollo del módulo de autenticación}

En esta sección se describe el desarrollo del módulo de autenticación, el cual se encarga de gestionar el acceso seguro de los usuarios. Para ello, se implementaron mecanismos basados en tokens JWT, control de acceso por roles y la definición de APIs que permiten la comunicación entre los distintos componentes de la aplicación web.

\subsubsection{Implementación del token JWT}

Para la gestión de autenticación en el sistema, se implementó un mecanismo basado en JSON Web Tokens (JWT), el cual permite validar la identidad del usuario de forma segura y sin necesidad de mantener sesiones en el servidor.

El proceso de autenticación inicia cuando el usuario envía sus credenciales a través del endpoint de inicio de sesión. Estas credenciales son verificadas contra la información almacenada en la base de datos, validando tanto la existencia del usuario como su estatus y la coincidencia de la contraseña cifrada.

Una vez que las credenciales son validadas correctamente, se genera un token JWT que contiene información relevante del usuario, como su identificador, nombre y tipo de usuario. Este token incluye además un tiempo de expiración, con el objetivo de limitar su validez y mejorar la seguridad del sistema.

El token generado es enviado al cliente, quien deberá incluirlo en las solicitudes posteriores para acceder a los recursos protegidos del sistema.

El proceso de generación del token se realiza mediante una función que recibe los datos del usuario y construye la carga útil (payload) del token. Posteriormente, se agrega el campo de expiración y se codifica utilizando una clave secreta y un algoritmo de cifrado.

A continuación, se muestra el fragmento de código utilizado para la generación del token JWT:

\begin{lstlisting}[language=Python, caption={Generación del token JWT}, label={lst:jwt_token}]
    def create_access_token(data: dict):
        to_encode = data.copy()
        expire = datetime.utcnow() + timedelta(minutes=ACCESS_TOKEN_EXPIRE_MINUTES)
        to_encode.update({"exp": expire})
        encoded_jwt = jwt.encode(to_encode, SECRET_KEY, algorithm=ALGORITHM)
        return encoded_jwt
\end{lstlisting}


\subsubsection{Implementación del control de acceso por rol}

El sistema implementa un mecanismo de control de acceso basado en roles (RBAC, por sus siglas en inglés), con el objetivo de restringir el acceso a los recursos y funcionalidades según el tipo de usuario autenticado.

Cada usuario posee un rol definido dentro del sistema, el cual se encuentra almacenado en la base de datos y es incluido dentro del token JWT generado durante el proceso de autenticación. Este rol es utilizado posteriormente para validar los permisos de acceso a los diferentes endpoints del sistema.

Durante la implementación, se definieron distintos tipos de usuario, tales como alumno y tutor, cada uno con permisos específicos. Por ejemplo, los alumnos pueden realizar actividades relacionadas con la captura y envío de ejercicios, mientras que los tutores tienen acceso a funciones de revisión, asignación y monitoreo del desempeño.

Para garantizar el cumplimiento de estas restricciones, se implementaron validaciones en los endpoints del sistema, donde se verifica el rol del usuario antes de permitir la ejecución de una operación. En caso de que el usuario no cuente con los permisos necesarios, el sistema responde con un error de autorización.

Este enfoque permite mantener la seguridad del sistema, evitando accesos no autorizados y asegurando que cada usuario interactúe únicamente con las funcionalidades correspondientes a su rol.

El rol del usuario es incluido dentro del payload del token JWT durante el proceso de autenticación, permitiendo que sea utilizado en cada solicitud para validar los permisos de acceso sin necesidad de consultar constantemente la base de datos.

A continuación, se muestra un ejemplo de validación de rol en un endpoint del sistema:

\begin{lstlisting}[language=Python, caption={Validación de acceso por rol}]
    if current_user["tipo_usuario"] != "tutor":
        raise HTTPException(
            status_code=403,
            detail="Acceso no autorizado"
        )
    \end{lstlisting}

\subsubsection{Implementación de APIs}

El endpoint de inicio de sesión se encarga de validar las credenciales del usuario y generar el token correspondiente. Este proceso incluye la verificación del estatus del usuario, así como la validación de la contraseña cifrada.

Una vez autenticado, se construye el payload del token con la información del usuario y se genera el JWT que será utilizado para futuras solicitudes.

El siguiente fragmento muestra la implementación del endpoint de autenticación:

\begin{lstlisting}[language=Python, caption={Endpoint de autenticación y generación de JWT}, label={lst:jwt_login}]
    @router.post("/login", response_model=LoginResponse)
    async def login(credentials: LoginRequest):
        conn = connect_to_database()
        if not conn:
            raise HTTPException(
                status_code=status.HTTP_500_INTERNAL_SERVER_ERROR, 
                detail="Error de conexión a la base de datos"
            )
    
        cursor = conn.cursor()
    
        try:
            query = """
                SELECT id_usuario, contrasena_cifrada, nombre, tipo_usuario, id_estatus 
                FROM Usuario
                WHERE username = ?
            """
            cursor.execute(query, credentials.username)
            user = cursor.fetchone()
    
            if not user:
                raise HTTPException(
                    status_code=status.HTTP_401_UNAUTHORIZED,
                    detail="Usuario o contraseña incorrectos"
                )
    
            id_usuario, contrasena_cifrada, nombre, tipo_usuario, id_estatus = user
    
            if id_estatus != 1:
                raise HTTPException(
                    status_code=status.HTTP_403_FORBIDDEN,
                    detail="El usuario no se encuentra activo"
                )
    
            if not verify_password(credentials.password, contrasena_cifrada):
                raise HTTPException(
                    status_code=status.HTTP_401_UNAUTHORIZED,
                    detail="Usuario o contraseña incorrectos"
                )
    
            token_payload = {
                "id_usuario": id_usuario,
                "nombre": nombre,
                "tipo_usuario": tipo_usuario
            }
    
            token = create_access_token(token_payload)
    
            return {
                "message": "Inicio de sesión exitoso",
                "token": token
            }
    
        except HTTPException:
            raise
        except Exception as e:
            raise HTTPException(
                status_code=status.HTTP_500_INTERNAL_SERVER_ERROR, 
                detail=f"Error: {str(e)}"
            )
        finally:
            cursor.close()
            conn.close()
\end{lstlisting}
\subsection{Desarrollo del módulo de ejercicios}

En esta sección se describe el desarrollo del módulo de ejercicios, 
el cual gestiona la consulta del catálogo, la asignación de ejercicios 
a los alumnos y el monitoreo de su progreso. Se implementaron 
funcionalidades diferenciadas por rol: el tutor puede consultar el 
catálogo, asignar ejercicios y gestionar sus fechas límite, mientras 
que el alumno accede únicamente a los ejercicios que le han sido 
asignados.


\subsubsection{Implementación del catálogo}
 
El catálogo de ejercicios se implementó como una tabla de solo lectura en la base de datos, donde cada registro contiene atributos como título, descripción y tipo de ejercicio. Los tutores no pueden crear ni modificar ejercicios, únicamente consultarlos y seleccionarlos para asignarlos a sus alumnos.
 
El endpoint \texttt{GET api/ejercicios/tutor/\{id\_usuario\}} retorna los ejercicios activos asignados por el tutor autenticado. La consulta realiza un \texttt{JOIN} entre \texttt{Ejercicios\_Tutor} y \texttt{Ejercicios}, filtrando únicamente las asignaciones con \texttt{id\_estatus = 1}. Incorpora una validación de autorización que impide que un tutor consulte los ejercicios de otro, retornando un error \texttt{HTTP 403} cuando el \texttt{id\_usuario} del token no coincide con el parámetro de la ruta.
 
\begin{lstlisting}[language=SQL, caption={Consulta de ejercicios asignados por el tutor}]
SELECT et.id_ejercicio_tutor, e.id_ejercicio, 
       e.titulo, e.descripcion, e.tipo
FROM Ejercicios_Tutor et
JOIN Ejercicios e ON et.id_ejercicio = e.id_ejercicio
WHERE et.id_usuario = ? AND et.id_estatus = 1
\end{lstlisting}

\subsubsection{Implementación de asignación de ejercicios}
 
La asignación de ejercicios se implementó mediante el endpoint \texttt{POST /api/ejercicios/tutor}, el cual permite al tutor seleccionar un ejercicio del catálogo y asignarlo a sus alumnos con una fecha límite de entrega opcional. Antes de registrar la asignación, el endpoint verifica que el tutor no cuente ya con una asignación activa del mismo ejercicio, retornando un error \texttt{HTTP 400} en caso de duplicado. Esta validación complementa la restricción \texttt{UQ\_ET\_activo} definida en el Código~\ref{lst:restricciones_unicidad}, garantizando que el mensaje de error sea descriptivo para el cliente.
 
\begin{lstlisting}[language=SQL, caption={Validación e inserción de asignación de ejercicio}]

SELECT id_ejercicio_tutor FROM Ejercicios_Tutor
WHERE id_ejercicio = ? AND id_usuario = ? AND id_estatus = 1
 
INSERT INTO Ejercicios_Tutor 
    (id_ejercicio, id_usuario, id_estatus, fecha_desactivacion)
VALUES (?, ?, 1, ?)
\end{lstlisting}

\subsubsection{Implementación de ejercios para alumnos}

Los alumnos pueden ver los ejercicios pendientes, completados y vencidos, se implemetó esta función mediante los endpoints:

\begin{itemize}
    \item \texttt{GET api/ejercicios/alumno/\{id\_alumno\}/pendientes}
    \item \texttt{GET api/ejercicios/alumno/\{id\_alumno\}/completados}
    \item \texttt{GET api/ejercicios/alumno/\{id\_alumno\}/vencidos}
\end{itemize}
    

Los cuales permiten visualizar los ejercicios clasificados y tener mayor control en sus ejercicios, mediante las siguientes consultas SQL:

\begin{lstlisting}[language=SQL, caption={Consulta de ejercicios pendientes}]

    SELECT 
        et.id_ejercicio_tutor, e.id_ejercicio, e.titulo, 
        e.descripcion, e.tipo, e.contenido_base, et.fecha_desactivacion as fecha_fin
    FROM Ejercicios_Tutor et
    JOIN Ejercicios e ON et.id_ejercicio = e.id_ejercicio
    JOIN Usuario a ON a.id_tutor = et.id_usuario AND a.tipo_usuario = 'alumno'
    WHERE a.id_usuario = ?
        AND et.id_estatus = 1
        AND et.id_ejercicio_tutor NOT IN (
    SELECT id_ejercicio_tutor FROM Intentos WHERE id_usuario = ? )

\end{lstlisting}


\begin{lstlisting}[language=SQL, caption={Consulta de ejercicios completados}]

    SELECT *
    FROM vw_historial_alumno
    WHERE id_alumno = ?
        AND tipo_intento = 'Original'
    ORDER BY fecha_envio DESC

\end{lstlisting}

Para gestionar automáticamente los ejercicios vencidos, se implementó un proceso programado utilizando la librería \texttt{APScheduler}. Este proceso tiene como objetivo actualizar periódicamente el estado de las asignaciones cuya fecha límite ha expirado, evitando que dicha validación tenga que realizarse manualmente desde el cliente o mediante consultas complejas en tiempo real.

La funcionalidad se implementó mediante dos funciones principales: 

\begin{itemize}
    \item \texttt{cerrar\_ejercicios\_vencidos()}
    \item \texttt{iniciar\_scheduler()}.
\end{itemize}

La función \texttt{cerrar\_ejercicios\_vencidos()} establece una conexión con la base de datos y ejecuta una consulta \texttt{UPDATE} sobre la tabla \texttt{Ejercicios\_Tutor}. La consulta modifica el campo \texttt{id\_estatus} de las asignaciones activas (\texttt{id\_estatus = 1}) cuya fecha de desactivación ya expiró, cambiando su valor a \texttt{2}, correspondiente al estado de ejercicio vencido.

Por otra parte, la función \texttt{iniciar\_scheduler()} registra la tarea automática dentro del programador asíncrono, configurando su ejecución periódica cada hora. De esta manera, el sistema mantiene actualizados los estados de las asignaciones sin intervención manual.

\begin{lstlisting}[language=Python, caption={Proceso automático para cierre de ejercicios vencidos}]
scheduler = AsyncIOScheduler()

def cerrar_ejercicios_vencidos():
    conn = connect_to_database()
    if not conn:
        print("Cron: Error de conexion")
        return

    cursor = conn.cursor()

    try:
        cursor.execute("""
            UPDATE Ejercicios_Tutor
            SET id_estatus = 2
            WHERE id_estatus = 1
            AND fecha_desactivacion IS NOT NULL
            AND fecha_desactivacion < GETDATE()
        """)

        conn.commit()
        print(f"Cron: {cursor.rowcount} ejercicios cerrados")

    except Exception as e:
        print("Cron error:", e)

    finally:
        cursor.close()
        conn.close()

def iniciar_scheduler():
    scheduler.add_job(
        cerrar_ejercicios_vencidos,
        'interval',
        hours=1
    )
    scheduler.start()
\end{lstlisting}

La implementación de este proceso automático permite optimizar las consultas realizadas por el sistema, mantener consistencia en los estados de las asignaciones y garantizar que los alumnos visualicen correctamente los ejercicios vencidos en tiempo real.
\subsection{Desarrollo del módulo de captura de escritura}

El módulo de captura de escritura tiene como objetivo permitir que el alumno registre evidencia escrita de los ejercicios realizados para su posterior análisis mediante los módulos OCR, NLP y evaluación caligráfica. Este módulo fue desarrollado utilizando tecnologías web del lado del cliente, integrando herramientas de captura manual, carga de imágenes y acceso a cámara desde el navegador.

\subsubsection{Implementación del lienzo digital}

El módulo de lienzo digital fue desarrollado utilizando React y la API Canvas de HTML5, permitiendo al alumno realizar ejercicios de escritura directamente desde el navegador. La implementación incluye herramientas de lápiz y borrador, así como funcionalidades para deshacer acciones, limpiar el área de trabajo y exportar el dibujo realizado.

Para la captura de trazos se implementaron eventos de mouse y dispositivos táctiles, permitiendo compatibilidad tanto en computadoras como en dispositivos móviles. El contexto del canvas se configura dinámicamente dependiendo de la herramienta seleccionada, ajustando propiedades como color, grosor y tipo de operación gráfica.

El siguiente fragmento muestra la lógica principal utilizada para iniciar y realizar el dibujo dentro del lienzo:

\begin{lstlisting}[language=JavaScript, caption={Implementación básica del lienzo digital}]
const iniciarDibujo = (e) => {
    const { x, y } = obtenerCoordenadas(e);

    if (esBorrador) {
        contextRef.current.globalCompositeOperation = 'destination-out';
        contextRef.current.lineWidth = grosurGoma;
    } else {
        contextRef.current.globalCompositeOperation = 'source-over';
        contextRef.current.strokeStyle = 'black';
        contextRef.current.lineWidth = grosurLapiz;
    }

    contextRef.current.beginPath();
    contextRef.current.moveTo(x, y);
    setIsDrawing(true);
};

const dibujar = (e) => {
    if (!isDrawing) return;

    const { x, y } = obtenerCoordenadas(e);

    contextRef.current.lineTo(x, y);
    contextRef.current.stroke();
};
\end{lstlisting}

Adicionalmente, se implementó un mecanismo de historial basado en imágenes serializadas del canvas, permitiendo restaurar estados anteriores mediante la funcionalidad de deshacer.

Una vez finalizado el ejercicio, el contenido del lienzo se convierte en una imagen codificada en Base64 y se envía al backend mediante una solicitud HTTP autenticada.

\begin{lstlisting}[language=JavaScript, caption={Exportación y envío del dibujo al backend}]
const imagenDataUrl = canvasRef.current.toDataURL('image/png');

await fetch('/api/intentos/registrar', {
    method: 'POST',
    headers: {
        'Content-Type': 'application/json',
        'Authorization': `Bearer ${token}`
    },
    body: JSON.stringify({
        id_ejercicio_tutor: idEjercicioTutor,
        imagen_codificada: imagenDataUrl
    })
});
\end{lstlisting}

\subsubsection{Implementación de carga de imagen y cámara}

El módulo de captura de escritura fue desarrollado en React, permitiendo al alumno registrar evidencias mediante la cámara del dispositivo o la carga de imágenes desde el almacenamiento local. Su objetivo es facilitar la adquisición de muestras de escritura para posteriormente ser procesadas por los módulos OCR y NLP.

La implementación valida el tipo y tamaño de los archivos recibidos, aceptando únicamente imágenes con un límite máximo de 10 MB. Una vez seleccionada la imagen, esta es redimensionada y comprimida utilizando un elemento \texttt{canvas}, optimizando así el envío y almacenamiento de la información.

\begin{lstlisting}[language=JavaScript, caption={Procesamiento y validación de imágenes}]
const procesarArchivo = (file) => {
    if (!file.type.startsWith('image/')) return;
    if (file.size > 10 * 1024 * 1024) return;

    const reader = new FileReader();

    reader.onloadend = () => {
        const img = new Image();

        img.onload = () => {
            const canvas = document.createElement('canvas');
            canvas.width = 1200;
            canvas.height = 800;

            canvas.getContext('2d')
                  .drawImage(img, 0, 0, 1200, 800);

            setImagenPreview(
                canvas.toDataURL('image/jpeg', 0.8)
            );
        };

        img.src = reader.result;
    };

    reader.readAsDataURL(file);
};
\end{lstlisting}

Para la captura desde cámara web se utilizó la API \texttt{navigator.mediaDevices.getUserMedia()}, permitiendo acceder al dispositivo de video y obtener fotografías directamente desde la aplicación.

\begin{lstlisting}[language=JavaScript, caption={Inicialización de la cámara web}]
const iniciarCamaraWeb = async () => {
    const mediaStream =
        await navigator.mediaDevices.getUserMedia({
            video: {
                facingMode: 'environment'
            }
        });

    setStream(mediaStream);
    setMostrarCamaraWeb(true);
};
\end{lstlisting}

Asimismo, se incorporaron herramientas básicas de edición como rotación y recorte de imagen, con el propósito de mejorar la calidad de las capturas antes de ser enviadas al servidor.

\subsubsection{Implementación de procesamiento de imagen}

Antes de enviar las evidencias de escritura al servidor, las imágenes pasan por una etapa de preprocesamiento cuyo objetivo es optimizar su calidad y reducir el tamaño de almacenamiento y transmisión. Este proceso mejora la compatibilidad con los módulos OCR y NLP encargados del análisis posterior.

El preprocesamiento implementado contempla:

\begin{itemize}
    \item Validación del tipo y tamaño del archivo.
    \item Redimensionamiento de la imagen para mantener dimensiones uniformes.
    \item Compresión en formato JPEG para reducir el peso del archivo.
    \item Conversión de la imagen a formato Base64 para su envío mediante HTTP.
\end{itemize}

El redimensionamiento y compresión se realizaron utilizando un elemento \texttt{canvas}, sobre el cual se dibuja la imagen original antes de exportarla nuevamente.

\begin{lstlisting}[language=JavaScript, caption={Preprocesamiento de imagen antes del envío}]
const canvas = document.createElement('canvas');

canvas.width = width;
canvas.height = height;

const ctx = canvas.getContext('2d');

ctx.drawImage(img, 0, 0, width, height);

const imagenProcesada =
    canvas.toDataURL('image/jpeg', 0.8);

setImagenPreview(imagenProcesada);
\end{lstlisting}

Finalmente, la imagen procesada es enviada al backend mediante una solicitud HTTP para su almacenamiento y posterior análisis automático.

